\section{Theoretical Background} \label{sec:theoreticalbackground}

% %=========================================================


\subsection{Static Code Analysis and Bug Detection}
 
\hspace{0.6em} \textbf{\gls{sca} examines a program without executing it or supplying specific runtime inputs}. On the other side there is Dynamic Code Analysis (or Code Dynamic Analysis) that evaluates a software program while it is actively running to observe its real time behaviour, performance and security. Static analysis inspects source code or compiled representations to identify potential defects before they cause failures in runtime. One of the firsts \gls{sca} tools actually dates back on 1977, with the paper of \textit{Johnson Stephen C.} "\textit{Lint, a C program checker}".

\hspace{0.6em} A newer tool was studied in 2008 by \textit{Ayewah et al.} in their paper "\textit{Using static analysis to find bugs}" \cite{paperreference1_ayewah2008using}, called \textbf{FindBugs}. This tool analyzes compiled Java bytecode using detectors designed to recognize common programming mistakes. Examples of detectable defects include potential null pointer dereferences, infinite recursion, ignored return values and logically inconsistent comparisons.

\hspace{0.6em} A reported defect is a \textbf{true positive when it corresponds to an actual programming error} and a \textbf{false positive when the reported condition is not a genuine defect}. Traditional analyzers cannot always determine whether a suspicious execution path is feasible. Consequently, they may report false positives or miss genuine defects. At the respective time, \textit{FindBugs} deliberately focused on defects that can be detected efficiently rather than guaranteeing detection of every possible error.

\hspace{0.6em} A paper from 2010 titled "\textit{A few billion lines of code later...}" by \textit{Bessey et al.} \cite{paperreference2_bessey2010few} took the topic and applied it on a larger scale. They demonstrate that on an industrial level the practical static analysis must be adapted to handle large codebases, complex build environments and diverse programming constructs. They also state that excessive false positives undermine developer trust and increase the effort required to review warnings.


% %=========================================================

\subsection{Machine Learning for Static Analysis}

\hspace{0.6em} Modern solution implies \textbf{supervised learning} by using labeled examples to train a \gls{ml} model. \textbf{The model learns from features such as warning types, surrounding code and program structure} to classify new warnings. \textbf{Code can be represented using numerical features, token sequences or graphs}. The surrounding code is important because a warning may depend on conditions or statements outside the reported line \cite{paper1_kharkarlearning2022,paper2_hegedHus2022static,paper3_yerramreddy2023empirical}.

\hspace{0.6em} A typical \gls{ml} assisted \gls{sca} process is:
\[
\underbrace{\text{Source Code}
\rightarrow \text{Static Analyzer}
\rightarrow \text{Warnings}}_{\gls{sca}}
\rightarrow
\underbrace{\text{ML Classifier}
\rightarrow \text{Warning Classification}}_{\gls{ml}}
\]

\hspace{0.6em} Unlike direct only \gls{ml} based bug detection, this approach uses \gls{ml} to filter warnings already generated by a static analyzer. \textbf{It can reduce false positives and help developers focus on genuine defects, but it cannot detect bugs that the original analyzer missed} \cite{paper1_kharkarlearning2022}.


\subsection{Evaluation Metrics}

\hspace{0.6em} \gls{ml} warning classifiers can be evaluated using precision, recall, F1-score and accuracy \cite{paper2_hegedHus2022static,paper3_yerramreddy2023empirical}. Considering real defects as the positive class, $TP$ represents correctly retained defects, $TN$ correctly rejected false warnings, $FP$ incorrect warnings retained and $FN$ real defects mistakenly filtered out: \newline

$\mathrm{Precision} = \frac{TP}{TP+FP},\quad \mathrm{Recall} = \frac{TP}{TP+FN},\quad \mathrm{F1} = 2\cdot\frac{\mathrm{Precision}\cdot\mathrm{Recall}}{\mathrm{Precision}+\mathrm{Recall}},\quad \mathrm{Accuracy} = \frac{TP+TN}{TP+TN+FP+FN}$ \newline

\hspace{0.6em} Precision shows how many retained warnings are real defects, while recall shows how many real defects are kept. $F1$ score balances these two measures and accuracy shows the proportion of correct classifications. High precision reduces unnecessary developer work, but recall must also be considered to avoid removing genuine warnings \cite{paper1_kharkarlearning2022,paper3_yerramreddy2023empirical}.


% %=========================================================
